\subsubsection{Conditional Language Models}

In the previous subsection, we defined the goal of sequence-to-sequence learning as modeling the conditional probability:
\begin{equation*}
    P\left(y_1, y_2, \dots, y_n \mid x_1, x_2, \dots, x_m\right)
\end{equation*}
However, directly \textbf{modeling this probability is intractable}. The output sequence is a structured object whose number of possible configurations grows exponentially with its length. For a vocabulary (i.e., the set of all possible tokens) of size $\left|\mathcal{V}\right|$, there are $\left|\mathcal{V}\right|^{n}$ possible sequences of length $n$. Therefore, assigning a probability to each possible sequence as a whole is \textbf{computationally infeasible}.

\highspace
To make this problem manageable, we use the \textbf{chain rule of probability} to decompose the joint probability of the sequence into a product of conditional probabilities:
\begin{equation}
    P\left(y_1, y_2, \dots, y_n \mid x\right)
    =
    \prod_{t=1}^{n}
    P\left(y_t \mid y_1, y_2, \dots, y_{t-1}, x\right)
\end{equation}
This formulation is known as a \definition{Conditional Language Model}, since it extends the idea of a standard language model by conditioning on an external input $x$.

\highspace
\begin{definitionbox}[: Language Model]
    A \definition{Language Model (LM)} is a probabilistic model that assigns a probability to a sequence of tokens:
    \begin{equation*}
        y = \left(y_1, y_2, \dots, y_n\right)
    \end{equation*}
    By modeling the joint distribution:
    \begin{equation*}
        P\left(y_1, y_2, \dots, y_n\right)
    \end{equation*}

    \highspace
    Using the chain rule of probability, this distribution is factorized as:
    \begin{equation}
        P\left(y_1, y_2, \dots, y_n\right)
        =
        \prod_{t=1}^{n}
        P\left(y_t \mid y_1, y_2, \dots, y_{t-1}\right)
    \end{equation}

    \highspace
    Thus, a language model predicts each token $y_t$ based on the sequence of previous tokens:
    \begin{equation*}
        y_1 \rightarrow y_2 \rightarrow \dots \rightarrow y_n
    \end{equation*}
    In simple terms, a language model learns to predict the next word in a sequence given the words that came before it. For example, given the input sequence ``\emph{The cat is on the}'', a language model might predict that the next word is likely to be ``\emph{mat}'' or ``\emph{roof}'' based on the context provided by the previous words. It is important to note that \hl{a language model does not have access to any external input; it only relies on the sequence of tokens it has seen so far.} Furthermore, a \textbf{language model does not understand the meaning of the words}; it simply learns statistical patterns in the data to make predictions. For instance, it might learn that the word ``\emph{cat}'' is often followed by ``\emph{is}'' and then by ``\emph{on}'' in many sentences, without any understanding of what a cat is or what being on something means.
\end{definitionbox}

\begin{flushleft}
    \textcolor{Green3}{\faIcon{question-circle} \textbf{What is a token?}}
\end{flushleft}
A \definition{Token} is a basic unit of text that a language model processes. It can be a word, a subword, or even a character, depending on the tokenization scheme used. More generally, a token represents the basic unit generated by the model at each timestep. For example, in the sentence ``\emph{The cat is on the mat}'', the tokens could be:
\begin{itemize}
    \item Words: ``The'', ``cat'', ``is'', ``on'', ``the'', ``mat''.
    \item Subwords: ``Th'', ``e'', ``cat'', ``is'', ``on'', ``the'', ``mat''.
    \item Characters: ``T'', ``h'', ``e'', `` '', ``c'', ``a'', ``t'', `` '', ``i'', ``s'', `` '', ``o'', ``n'', `` '', ``t'', ``h'', ``e'', `` '', ``m'', ``a'', ``t''.
\end{itemize}
The choice of tokenization can affect the performance of the language model, as it determines the granularity of the input and output sequences. For example, using subword tokenization can help the model handle rare words and morphological variations more effectively than word-level tokenization.

\highspace
\begin{definitionbox}[: Conditional Language Model]
    A \textbf{language model} assigns a probability to a sequence of tokens:
    \begin{equation*}
        P\left(y_1, y_2, \dots, y_n\right)
        =
        \prod_{t=1}^{n}
        P\left(y_t \mid y_1, \dots, y_{t-1}\right)
    \end{equation*}
    That is, each token is predicted based on the previous ones.

    \highspace
    A \definition{Conditional Language Model} extends this formulation by conditioning on an additional input $x$:
    \begin{equation*}
        P\left(y_1, y_2, \dots, y_n \mid x\right)
        =
        \prod_{t=1}^{n}
        P\left(y_t \mid y_1, \dots, y_{t-1}, x\right)
    \end{equation*}
    Thus, each token is generated based on both:
    \begin{itemize}
        \item The previously generated tokens;
        \item The external input $x$.
    \end{itemize}
    In other words, a conditional language model \textbf{learns to generate a sequence of tokens that is not only coherent with itself but also relevant to the given input} $x$. For example, if $x$ is an English sentence and $y$ is its French translation, the model learns to generate a French sentence that is both grammatically correct and semantically related to the English input. This makes conditional language models particularly powerful for tasks like machine translation, where the output must be closely tied to the input.
\end{definitionbox}

\noindent
\textcolor{Green3}{\faIcon[regular]{lightbulb} \textbf{Key idea.}} The output sequence is generated \textbf{one element at a time}. At each timestep $t$, the model predicts the next token $y_t$ based on:
\begin{itemize}
    \item The input sequence $x$;
    \item The previously generated outputs $y_1, \dots, y_{t-1}$.
\end{itemize}
This transforms sequence generation into a step-by-step process, where the model maintains a memory of the past through its hidden state and uses it to inform future predictions.

\highspace
\textcolor{Green3}{\faIcon{question-circle} \textbf{Why is this important?}} This factorization makes the problem tractable: instead of predicting an entire sequence at once, the model learns to predict one token at a time, while maintaining a memory of the past through its hidden state.

\highspace
\textcolor{Green3}{\faIcon{info-circle} \textbf{Remark.}} For simplicity, the dependence on the model parameters $\theta$ is often omitted in the notation. However, all probabilities should be understood as being parameterized, i.e., $P\left(y_t \mid y_{<t}, x\right) = P\left(y_t \mid y_{<t}, x; \theta\right)$.